\chapter{Что на самом деле говорит дофамин}
% полная версия: chapters/07_dofaalt.tex

В прошлой главе дофамин был простым: одно число, «лучше или хуже, чем ожидалось». Это первое приближение, и оно хорошо работает. Но опыты последних лет показали, что по этому «проводу» идёт больше, чем одно число. Эта глава --- о том, где картина сложнее и где учёные пока спорят.

\section*{Оптимисты и пессимисты}

Дофаминовые нейроны не одинаковы. Одни сильнее реагируют на приятные сюрпризы, чем на неприятные, --- это оптимисты. Другие наоборот --- пессимисты. Если каждый нейрон учится по своему правилу, оптимист выучит, какой будет награда в лучшем случае, а пессимист --- в худшем. Вместе они записывают не одно среднее ожидание, а весь разброс возможных исходов, как букмекер, который знает не только фаворита, но и шансы каждого. Согласие с опытом у этой идеи есть; зачем мозгу весь разброс --- пока гипотеза.

\pencilfig{7}{\pencilart{7}}{Каждый дофаминовый нейрон держит свою точку на кривой возможных наград: вместе они помнят всю кривую.}

\section*{Не только «лучше», но и «важно»}

Часть дофаминовых нейронов вспыхивает и на неприятные, но важные события: громкий звук, опасность. Похоже, что по дофаминовой сети идёт не только знак сюрприза, но и его значимость --- сигнал «обрати внимание».

\section*{Два сигнала на одном проводе}

Быстрые всплески учат. А медленный средний уровень дофамина, меняющийся за минуты, похоже, говорит о другом: насколько богата сейчас обстановка. Если награды вокруг много, время дорого --- и животное действует быстрее. Опыты подтверждают, что уровень дофамина связан с энергичностью действий. Инженер назвал бы это разделением сигналов по частоте: высокие частоты несут одно, низкие --- другое.

\section*{Нарастание у цели}

Когда животное приближается к награде, дофамин часто не вспыхивает, а плавно растёт. Одно объяснение: это само ожидание, сигнал «цель близко, поднажми». Другое сохраняет идею сюрприза: издалека животное оценивает обстановку размыто, по мере приближения всё точнее, и каждое уточнение --- маленький приятный сюрприз. Изящный опыт в виртуальном коридоре, где животное мгновенно переносили ближе к цели, говорит в пользу второго: дофамин отвечал всплеском, как и положено сюрпризу.

\section*{Взгляд назад}

Есть и смелая альтернатива: дофамин отвечает не на вопрос «что будет?», а на вопрос «что было причиной награды?» --- смотрит в прошлое, а не в будущее. Опыты, где две теории предсказывают разное, пока противоречат друг другу. Это живой научный спор.

\section*{Вектор вместо числа}

Наконец, дофамин реагирует и на смену \emph{вкуса} награды, когда её ценность не изменилась, и помогает учиться без всякой награды. Похоже, это не одна ошибка, а целый пучок ошибок --- по одной на каждую черту того, что произошло.

Вывод главы: большинство новых теорий не отменяет простой картины, а расширяет её. Дофамин --- не один провод, а небольшой жгут из нескольких. Только «взгляд назад» спорит с ней всерьёз.

\fullversion{7}
